% =====================================================================
%  Chapter 12: Temperature and chemical potential
%  Gattringer & Lang, "Quantum Chromodynamics on the Lattice" (LNP 788)
%  Reconstructed from extract/ch12.txt; equations auto-numbered.
% =====================================================================

\chapter{Temperature and chemical potential}

Since the big bang the universe has cooled down substantially. Still,
understanding QCD at high temperature and at high matter density is important
for various reasons. One of them is obviously a better understanding of quark
and gluon matter shortly after the big bang and of the condensation of hadrons
in the cooling process. Many objects in the universe, like neutron stars,
supposedly have high enough density or temperature to expect that hadronic
matter behaves differently compared to usual atoms. On earth very energetic
collisions of heavy ions provide also possibilities to study matter under
extreme conditions. Thus studying QCD at high temperature and density from
first principles is a challenging task for the lattice approach.

Intuitively one may expect that at a density higher than in the nucleons and at
a temperature above the pion mass the hadrons loose their individuality as
bound states and new structural phases emerge. One anticipates phase
transitions and the lattice formulation is a natural setting for the study of
such collective phenomena. The studies should reveal the existence and the kind
of transitions to different states of matter and their properties, as well as
the behavior of the possible other states of matter in an unusual environment,
thus helping us to understand experiments as well as astrophysical
observations. Reviews for progress in lattice thermodynamics are, e.g., found
in [1--5].

In today's lattice approach we always study a system in equilibrium and not the
evolution and decay dynamics of the system. This challenge is left for the
future.

\section{Introduction of temperature}

In Sect.~1.4.3 we have discussed the relation of a Euclidean quantum field
theory to the partition function (1.13) of a classical statistical mechanics
system in a heat bath with temperature $T$. For the quantum mechanical case the
partition function is given by
\begin{equation}
  Z(T) = \tr\!\left[e^{-H/(k_B T)}\right] = \tr\!\left[e^{-\beta H}\right],
\end{equation}
where $\hat{H}$ is the Hamiltonian operator and $\beta$ now denotes the inverse
temperature $\beta = 1/(k_B T)$, with the Boltzmann constant $k_B$. This is the
standard notation, which, unfortunately, often leads to confusing $\beta$ with
the inverse gauge coupling. Unless stated otherwise, in this chapter $\beta$
exclusively denotes the inverse temperature. In our subsequent discussion we
follow the usual choice of units setting $k_B = 1$. Then the temperature $T$ is
given in units of energy or mass and $\beta = 1/T$.

Due to the trace in (12.1) we are restricted to fields that are periodic
(bosons) or anti-periodic (fermions) in time. Like in Chap.~1 the partition
function can be transformed to a path integral over such field configurations.
However, contrary to our discussion there, we now do not assume that the time
extent becomes infinite. Instead we obtain the functional integral ($\Phi$
being some generic field)
\begin{equation}
  Z(T) = \int \mathcal{D}[\Phi]\,e^{-S_E[\Phi]},
\end{equation}
where the integration is over fields, (anti-)periodic in the finite time
direction. The Euclidean action results from an integral over all space but
finite time extent:
\begin{equation}
  S_E[\Phi] = \int_0^{\beta}\!\dd t \int_{\mathbb{R}^3}\!\dd^3 x\,
    \mathcal{L}_E\big(\Phi(t,x),\partial_\mu\Phi(t,x)\big).
\end{equation}

The measure $\mathcal{D}[\Phi]$ and the action $S_E[\Phi]$ are discretized on a
lattice as usual. In the discretization we find that only our point of view
changes. Until now, working at ``zero temperature'' we were interested in
results in the infinite space--time volume limit and we discussed deviations as
finite size effects. Space and time extent were considered as much larger than
the largest correlation length in the system -- that of the pion. Now space is
still considered in that limit, but the physical extent of time is limited to
$\beta$. For a finite lattice the space extent is $aN$ and the time extent is
$aN_T$. We thus have
\begin{equation}
  \beta = a\,N_T = \frac{1}{T}
\end{equation}
and find that the limit $\beta\to\infty$ corresponds to $T\to 0$. We interpret
this as a system with finite spatial volume and fixed temperature $T$. The
continuum limit of such a system corresponds to $a\to 0$ while holding $aN$ and
$aN_T$ fixed. Finite volume effects become smaller when $N/N_T$ becomes larger.

Arguing again with physical intuition we may expect to observe finite
temperature effects when the correlation length of the pion $1/M_\pi$ reaches
the time extent $aN_T$, and it has been found that $T$ is indeed not far from
the value $M_\pi$.

The fact that the temporal extent $\beta$ is kept finite has important physical
consequences. If Fourier transformation is applied to the finite time direction
one finds that only discrete energy levels with spacing $\Delta = 2\pi/\beta$
are allowed. Integer multiples of this value $\Delta$ are often referred to as
Matsubara frequencies. Due to the lattice structure the Matsubara frequencies
are not only discrete but also limited to the range $(-\pi/a,\pi/a]$, taking
for boson states values $k\Delta$ (for $-N_T/2 + 1 \le k \le N_T/2$) with
$\Delta = 2\pi T = 2\pi/(aN_T)$. For fermions, which obey anti-periodic
boundary conditions in time direction, the values are shifted by $\Delta/2$
(cf.\ Appendix A) and the smallest quark Matsubara mode thus has energy
$\Delta/2$.

\subsection{Analysis of pure gauge theory}

Let us first discuss pure gauge theory at finite temperature. For this case we
introduced in Sect.~3.3.5 an important observable, the Polyakov loop. It is
defined as the trace of the ordered product of gauge link variables in time
direction, closed due to periodicity:
\begin{equation}
  P(m) = \tr\!\left[\prod_{j=0}^{N_T-1} U_4(m,j)\right].
\end{equation}

The correlator $\ev{P(m)\,P(n)^\dagger}$ is related to the potential between a
pair of a static quark and antiquark (compare (3.61)), i.e., the free energy
$F_{\bar{q}q}$ of such a pair at a given temperature:
\begin{equation}
  \ev{P(m)\,P(n)^\dagger} = e^{-a N_T\,F_{\bar{q}q}(a|m-n|)}
    = e^{-F_{\bar{q}q}(r)/T}.
\end{equation}

As usual the energy has to be normalized at some distance. For large distances
we expect factorization
\begin{equation}
  \lim_{a|m-n|\to\infty} \ev{P(m)\,P(n)^\dagger}
    = \ev{P(m)}\,\ev{P(n)^\dagger} = \left|\ev{P}\right|^2.
\end{equation}

Due to translational invariance the spatial position of the Polyakov loop is
irrelevant and we have therefore omitted it on the right-hand side of (12.7)
and instead use the spatial average
\begin{equation}
  P = \frac{1}{N^3}\sum_m P(m).
\end{equation}

We observe that for static potentials that grow indefinitely with separation
distance, such as the confining potential (3.62) at $\sigma > 0$,
$\left|\ev{P}\right|$ has to vanish. Therefore, we conclude
\begin{equation}
  \ev{P} = 0 \;\Leftrightarrow\; \text{confinement},
  \qquad
  \ev{P} \ne 0 \;\Leftrightarrow\; \text{no confinement}.
\end{equation}

At low temperatures QCD is confining. As the temperature is increased, pure
gauge theory undergoes a phase transition [6, 7] at a critical temperature
$T_c$ of about 270~MeV. There the system deconfines and $\ev{P}$ acquires a
nonvanishing value (cf., the left-hand side of Fig.~12.1).

An individual Polyakov loop $P(m)$ is nothing but a Wilson loop at fixed
spatial position extending in the temporal direction and closed due to
periodicity. We may interpret its expectation value as the probability to
observe a single static charge. With
\begin{equation}
  \left|\ev{P}\right| \sim e^{-F_q/T}
\end{equation}
we relate it to the free energy of a single color charge. From (12.9) we
conclude that $F_q \to \infty$ for confinement and $F_q$ finite in the
deconfined phase.

The deconfinement transition in quenched QCD has another interesting
interpretation in terms of the so-called center symmetry or $Z_3$ symmetry.
From the definition of the Polyakov loop in (3.60) it is obvious that it is a
gauge-invariant quantity and cannot be made unity by a gauge transformation.
However, on a finite lattice any given gauge configuration has the same
statistical weight as the one obtained by a center transformation. This
transformation consists of multiplying all temporal links in a given time slice
$n_4 = t_0$ with the same element $z$ of the center group $Z_3$ of gauge group
$\mathrm{SU}(3)$:
\begin{equation}
  U_4(n,t_0) \;\to\; z\,U_4(n,t_0).
\end{equation}

The center elements of $\mathrm{SU}(3)$ are the cubic roots
$(1, e^{2\pi i/3}, e^{-2\pi i/3})$. The reason for this symmetry is that the
gauge action is constructed from products of link variables along trivially
closed loops. Each trivially closed loop has as many elements in one direction
as in the opposite. Center elements commute with all group elements and
therefore cancel in the closed loop products. For example, a plaquette in the
$(\mu,4)$-plane at $n_4 = t_0$ transforms according to
\begin{equation}
\begin{split}
  &\tr\Big[U_\mu(n,t_0)\,U_4(n+\hat{\mu},t_0)\,U_\mu(n,t_0+1)^\dagger\,
    U_4(n,t_0)^\dagger\Big] \\
  &\quad\longrightarrow\;
  \tr\Big[U_\mu(n,t_0)\,z\,U_4(n+\hat{\mu},t_0)\,U_\mu(n,t_0+1)^\dagger\,
    U_4(n,t_0)^\dagger\,z^\dagger\Big] \\
  &\quad= \tr\Big[z\,z^\dagger\,U_\mu(n,t_0)\,U_4(n+\hat{\mu},t_0)\,
    U_\mu(n,t_0+1)^\dagger\,U_4(n,t_0)^\dagger\Big].
\end{split}
\end{equation}
Since $zz^\dagger = 1$, the gauge action is invariant. Plaquettes in spatial
planes or with $n_4 \ne t_0$ are not affected. This global symmetry is the
heart of the proof of confinement of static charges in the strong coupling
limit in Sect.~3.4.

The Polyakov loop, however, does not close in a topologically trivial way, but
instead winds once around the compact time direction. Thus it is not invariant
under transformation (12.11). Since we multiply the temporal links in only one
time slice with a center element $z$, the Polyakov loop transforms as
\begin{equation}
  P \;\to\; z\,P.
\end{equation}

Due to the symmetry of the action the expectation value of the Polyakov loop may
be written as
\begin{equation}
  \ev{P} = \frac{1}{3}\left(P + zP + z^2P\right)
    = \frac{1}{3}\left(1 + e^{i2\pi/3} + e^{-i2\pi/3}\right)\ev{P} = 0 .
\end{equation}
The right-hand side vanishes since the sum over the center elements is zero.
The analysis (12.14) breaks down, when the center symmetry is broken. In
particular this symmetry is broken spontaneously above the critical value of
the temperature. Thus the finite temperature transition of quenched QCD signals
a spontaneous breaking of the center symmetry [8]. The Polyakov loop (also
called thermal Wilson line) is an order parameter distinguishing between a
confinement phase, where free charges cannot be found, and a deconfinement
phase, where single charges are screened and may be observed.

Combining several of such charges leads to nonvanishing expectation values as
long as the charges combine to a color singlet. This allows the definition of,
e.g., the static potential between static charges. Also the evolution of flux
tubes (concentration of gluonic energy) in situations of two or more static
charges has been studied [9--13].

The center symmetry property resembles that of a three-state spin model with
spin variables $\in Z_3$ living on the sites of the 3D spatial lattice volume
[14, 15]. Such a model exhibits a phase transition at some value of its nearest
neighbor coupling -- but only for an infinite volume. In a finite volume the
symmetry cannot be broken in a strict sense. However, computer simulations at
coupling values, where one would expect symmetry breaking in the infinite
volume situation, may stay for long persistence periods in one of the three
sectors, before a tunneling to another sector occurs.

The definition and measurement of the order parameter are a delicate issue. In
spin systems one introduces an external field, lets the volume go to infinity
and, then the external field to zero. This order of limits is important, since
the spontaneous magnetization vanishes for any finite volume. For the gauge
theory we have a similar situation preventing the correct determination of
$\ev{P}$ in a finite volume. Technically, larger lattices lead to long periods
in one of the sectors and thus one may expect good approximate determinations
of this would-be order parameter. However, the systematic error due to
tunneling is not under control. In the updating procedure one can even
explicitly enforce the center symmetry by regularly multiplying all gauge links
in a time slice with a nontrivial center element. This does not change the
gauge action. Then it becomes obvious that $\ev{P}$ is bound to vanish for
sufficiently long runs. Instead of $\ev{P}$ one therefore often plots and
analyses $\ev{|P|}$, which agrees with $\left|\ev{P}\right|$ in the infinite
volume limit.

In order to simulate at finite temperature $1/(aN_T)$ one either changes $N_T$
or the gauge coupling and therefore the lattice spacing $a$. Changing $a$
allows for a continuous scan over $T$. In Fig.~12.1 we show the result of a
simulation of pure gauge theory on a lattice of size $16^3\times 4$ for a range
of coupling values such that $T$ is close to the transition point from
confinement to deconfinement. In the left-hand side plot we show $\ev{|P|}$ as
a function of $T$. The right-hand side is a scatter plot of $P$ for individual
configurations at a fixed value of $T > T_c$. In the simulation also the center
symmetry was updated and the values of $P$ concentrate in the directions of the
center elements.

\begin{figure}[htbp]
  \centering
  \includegraphics[width=0.9\textwidth]{fig121.png}
  \caption*{Fig. 12.1. \emph{L.h.s.: The expectation value $\ev{|P|}$ as
    function of the temperature $T$. In computing $T = 1/(N_T a) = 1/(4a)$ the
    lattice spacing has been determined using the static potential as discussed
    in Sect.~3.5, cf.\ (3.83) (for $r_0 = 0.5$~fm). The results are from a
    simulation on lattice size $16^3\times 4$ averaging over 100 measured
    configurations for each value of the gauge coupling (i.e., the
    temperature); for a realistic determination better statistics and an
    extrapolation to infinite volume should be performed. R.h.s.: Scatter plot
    of Polyakov loop results in the complex plane for 1500 individual gauge
    configurations. In the broken phase the values of $P$ concentrate around
    the three center phases. The data are from a run with $\ev{|P|} = 0.53$.}}
  \label{fig:121}
\end{figure}

The results in general will be affected by the finiteness of the spatial volume
and by scaling corrections. Comparing different parameter settings allows one
to study both.

Monte Carlo calculations for the quenched theory have been performed already
very early [16--18] and improved considerably since then [19, 20]. Careful
studies revealed that the phase transition is weakly first order. This is in
agreement with the analogy to the $Z_3$-spin model [14].

In the determination of the transition temperature the lattice spacing enters
and the scale has to be fixed, as discussed in Sect.~3.5. The currently
accepted value is $T_c \approx 270$~MeV for the quenched situation. Results from
different improved gauge actions, extrapolated to the continuum limit, agree
within a few percent.

In studies of the transition various standard tools of statistical physics have
been employed. An informative quantity is the Polyakov loop susceptibility
\begin{equation}
  \chi_P = N^3\left(\ev{P^2} - \ev{P}^2\right),
\end{equation}
which peaks at the transition. The value and position of this peak follow
finite size scaling laws [21] and lead to information on the infinite volume
values from finite volume measurements.

Typical computer simulations at finite $T$ should have spatial extents much
larger than the temporal size in order to reduce finite size effects. Thus
$N_T$ is relatively small, which limits the number of different Matsubara
frequencies. It also complicates attempts to determine exponential decay
properties in time direction of observables like particle propagators since at
large $T$, i.e., small $N_T$, they can be followed only over very few time
steps. For that reason one also has introduced different lattice spacings $a$
and $a_T$ in spatial and temporal directions, respectively, with $a > a_T$.
This may be obtained by choosing different gauge couplings multiplying the
time-directed plaquettes and the purely spatial plaquettes. Keeping the ratio
$\xi = a/a_T$ large allows one to have large $N_T$ while keeping the
temperature constant. This method introduced in [22, 23] is discussed in more
detail in [24].

\subsection{Switching on dynamical fermions}

Dynamical fermions break the $Z_3$ symmetry explicitly and this affects the
Polyakov loop as well. A simple way to see this comes from the hopping
expansion of the fermion determinant introduced in Sect.~5.3.2. The exponent in
(5.65) may be considered as an effective fermionic action. It is a sum over
closed loops, like the gauge action itself. However, these are now all possible
loops on the lattice, including the ones that wind nontrivially around the
compactified time direction. In particular this sum will also include Polyakov
loops and therefore terms proportional to
\begin{equation}
  \kappa^{N_T}\left(P + P^*\right).
\end{equation}

Due to the transformation property (12.13) such a term in the effective action
is no longer invariant under an overall multiplication of all gauge links in a
time slice with an element of the group's center and therefore the fermions
break the center symmetry. In full QCD one finds that the Polyakov loop settles
around real positive values. In some sense one can consider the fermions' role
like that of an external, symmetry-breaking magnetic field in a spin model.

Another, more intuitive, picture is that the dynamically generated vacuum loops
screen the individual charges, similar to the Debye screening of electric
charges in a polarizable medium. The potential computed from the correlation
function (12.6) will saturate at some distance and may be approximated by a
screened Coulomb potential $\propto \exp(-\mu|r|)/|r|$. Eventually the notion
``potential'' looses its meaning, since it becomes energetically more favorable
to produce a quark--antiquark pair from the vacuum, saturating the test
charges: the ``string breaks.'' The Polyakov loop and the Wilson loop are no
longer true-order parameters. Confinement now has to be defined differently,
namely through the spectrum of the asymptotic states of the system, i.e., the
color neutral hadrons. Based on these considerations an alternative definition
of the order parameter was given in [25, 26].

In our theoretical world we can change the quark mass parameter $m$
continuously. Let us first discuss the case of two flavors with equal mass
$m_{u,d}$ (cf.\ Fig.~12.2). If we let the mass approach infinity, the quarks
decouple and we are in the quenched situation with a first-order temperature
phase transition. Decreasing the quark mass this transition persists but the
latent heat decreases and the transition becomes weaker. Eventually the
first-order transition line in the $(g,m)$ plane ends with a second-order
point.

\begin{figure}[htbp]
  \centering
  \includegraphics[width=0.7\textwidth]{fig122.png}
  \caption*{Fig. 12.2. \emph{The plot elaborates the phase structure in the
    $(T,m)$-plane for the $\Nf = 2$ sub-manifold (front side of the 3D plot in
    Fig.~12.3).}}
  \label{fig:122}
\end{figure}

The regions neighboring to the confinement and the deconfinement domains are
therefore analytically connected. Such a crossover is not unusual. Just think
of water in its three phases: ice, liquid, and vapor. Depending on the pressure
and temperature there may be phase transitions but also continuous
crossover-type transitions.

However, we are not yet done. Decreasing the quark mass further we finally
arrive at the situation of massless quarks, where the action is chirally
invariant (if a suitable discretization is used). As we have discussed in
Chap.~7, the full theory breaks the chiral symmetry spontaneously at zero
temperature. The corresponding order parameter is the quark condensate
$\Sigma$ defined in (7.63), which, however, is an order parameter only for the
massless theory. As one increases $T$ the condensate shows a transition. Its
value vanishes at some critical temperature $T_{ch}$ and stays zero for higher
temperatures. Thus the condensate may be considered a valid order parameter
indicating the chiral temperature transition (cf.\ Fig.~12.2).

In Sect.~7.3.5 we have related the chiral condensate to the spectral density of
the Dirac operator eigenvalues near the origin via the Banks--Casher relation
(7.77). This formula describes the situation also above the transition: At
$T_{ch}$ a gap opens up in the spectrum and the chiral condensate vanishes.

A highly nontrivial observation is that confinement also seems to go away at
about the same temperature $T_{ch}$ which suggests a strong connection between
confinement and chiral symmetry breaking. This equality of $T_{deconf}$ and
$T_{ch}$ has been disputed recently [1, 2, 27]. Possible reasons for the
controversy may be related to the definition of the pseudo-critical transition
point(s) and to the fermion species and masses used in the simulations.

The value of $T_{ch}$ is not easy to determine, since it requires a simulation
at vanishing quark masses. Also, it depends on the number of quark species
considered. Typical values obtained by various simulations and extrapolations
range between $T_{ch} \approx 150$~MeV [27] and $T_{ch} \approx 190$~MeV [28].
This -- compared to pure gauge theory -- smaller value may be understood from
the observation that in the quenched case the only bound states are rather
heavy glueballs. Therefore, larger temperatures are needed to generate a
glueball gas dense enough to enable deconfinement. There is some caution
advisable, though, as mentioned above.

Figure 12.3 illustrates our present knowledge on the phase diagram in a space
of two mass-degenerate quarks with mass $m_{u,d}$ and a third quark with mass
$m_s$. It is presently uncertain, whether for the physical quark masses
(indicated by a dotted circle in the plot) the transition is located in the
crossover or in the first-order domain.

\begin{figure}[htbp]
  \centering
  \includegraphics[width=0.9\textwidth]{fig123.png}
  \caption*{Fig. 12.3. \emph{Sketch of the phase structure in the
    $(T,m_{u,d},m_s)$-space for two mass-degenerate flavors (mass $m_{u,d}$)
    and another third flavor ($m_s$) and vanishing quark number
    density/chemical potential $\mu = 0$. The skewed shaded area indicates the
    position of the crossover transition and, where the shading is dark, the
    region where the transition is of first order. At the boundaries of the
    first-order region one has a second-order phase transition. The boundary
    lying in the $m_{u,d} = 0$ plane may have a second-order part, joining the
    first-order domain at a tricritical point. The figure is based on
    information discussed and summarized in [29].}}
  \label{fig:123}
\end{figure}

Lattice results with two light and with or without a heavier strange quark near
their physical mass values find a rapid but smooth crossover from lower to
higher temperature, but no convincing evidence for a universal critical
behavior has been found. Working at smaller lattice spacing, closer to the
continuum limit, and with improved fermion actions with smaller cutoff effects
may help to further clarify that situation [30, 31].

The order of the chiral temperature transition is less established than that
for the pure gauge theory and depends on the number of quark flavors. Pisarski
and Wilczek [32] have studied an effective 3D model field theory for the order
parameter, assuming the underlying spontaneously broken chiral symmetry of QCD.
A renormalization group analysis then led to the suggestion that the chiral
transition is second order for $\Nf = 2$ but weakly first order for $\Nf = 3$.
In the latter case the transition again extends toward $m > 0$ with a
second-order endpoint. This endpoint and the endpoint of the transition line
starting at the pure gauge theory are connected by a so-called crossover curve.
This is a maximum of some observables (like the specific heat) but not the
locus of nonanalyticity.

Like for the Polyakov loop, also the derivative of the condensate, the chiral
susceptibility
\begin{equation}
  \chi_{ch} = \frac{\partial\Sigma}{\partial m}
\end{equation}
is a useful quantity to locate the transition point, and it also has definite
finite size scaling behavior. Although neither $\Sigma$ nor $\ev{P}$ are strict
order parameters for $0 < m < \infty$ their values can still be used to study
the position of the phase transition lines.

\subsection{Properties of QCD in the deconfinement phase}

Formally the partition function depends on several parameters: the volume, the
temperature, the quark masses, and the gauge coupling. One defines
thermodynamic quantities like the energy density $\varepsilon$ or the pressure
$p$ by derivatives of the free energy ($F = -\ln(Z)/\beta$):
\begin{equation}
  \varepsilon \equiv -\frac{1}{V}\frac{\partial\ln Z}{\partial\beta}
    = \frac{T^2}{V}\frac{\partial\ln Z}{\partial T},
  \qquad
  p \equiv T\,\frac{\partial\ln Z}{\partial V},
  \qquad
  \Sigma_i = \frac{\partial\ln Z}{\partial m_i}.
\end{equation}
Here $V$ denotes the spatial volume and, as usual, $\beta = 1/T$.

As already addressed, the confinement and the deconfinement regions are
analytically connected. Still we may expect quite different properties of
matter in these regions, although the transition may be a gradual one. For
large temperature the scale for the running coupling is set by $T$ and this
brings one to the small gauge coupling perturbative regime. Asymptotic freedom
then lets one expect a gas of weakly interacting quarks and gluons: a plasma.
This is why the high-temperature phase is usually called the quark-gluon plasma
phase [33].

A relativistic ideal gas of (noninteracting) QCD gluons and quarks at high
temperature follows the Stefan--Boltzmann law and has the pressure and energy
density [34],
\begin{equation}
  p = \frac{\pi^2}{45}\left(8 + \frac{21}{4}\Nf\right) T^4,
  \qquad
  \varepsilon = 3p \quad \text{(equation of state)}.
\end{equation}
Due to the larger number of degrees of freedom of quarks and gluons this is
much larger than what one expects from a gas of, say, the three light pions
$\pi^0$ and $\pi^\pm$. It is not clear how quickly hadrons dissociate and such
a behavior is obtained. Figure 12.4 demonstrates that this may happen only at
quite high values of $T$.

\begin{figure}[htbp]
  \centering
  \includegraphics[width=0.9\textwidth]{fig124.png}
  \caption*{Fig. 12.4. \emph{Pressure $p$ determined from simulations for pure
    SU(3) gauge theory and QCD with two or three light quarks or two light and
    one heavy quark (2+1) with $m_s \approx T$ [35] (figure from [36]). The
    arrows on the right-hand side of the plot indicate the corresponding
    infinite temperature limits from the Stefan--Boltzmann law (12.19). (From:
    Karsch [36]. Reprinted with permission from AIP).}}
  \label{fig:124}
\end{figure}

Due to the crossover nature of the transition the particle spectrum may change
gradually from the (for zero temperature well-known) spectrum of mesons and
baryons to some intermediate state before asymptotic behavior sets in and a
plasma of quarks and gluons defines the ground state. The restoration of chiral
symmetry above $T_{ch}$ may be realized by degenerate parity partners, but
still, e.g., mesonic bound states. The region presently accessible in heavy ion
experiments is typically close to the transition and the observations thus may
be dominated and obscured by the intermediate phase-changing properties.

As mentioned, because of the small Euclidean time extent (less than 1~fm) it is
quite problematic to use the standard approach for determining the particle
spectrum by studying correlation functions in time direction, although this has
been tried [37--39]. One can, however, analyze the spatial correlation
functions, which leads to so-called screening masses [40--43]. Integrating the
propagators over all space and time gives the thermal susceptibilities
involving a contact term, which has to be subtracted. The integration involves
a sum over all states in the corresponding quantum channel, though. Following
these through the transition indicates that the masses of $\pi(0^{-+})$ and the
isoscalar $f_0(0^{++})$ become degenerate quickly above $T_{ch}$, indicating a
restoration of SU(2) flavor symmetry. The mass of the isovector $a_0(0^{++})$
is still different from that of the pion, though, indicating persistent
$U_A(1)$ symmetry breaking.

\section{Introduction of the chemical potential}

The net number of baryons or quarks in the vacuum is zero, but it is not in
matter. In extreme situations, like heavy ion collisions or ultra-dense matter
in neutron stars, one has to consider the effect of nonvanishing quark number
density $n_q$. In this case we have to extend the partition function of the
grand canonical ensemble (12.1):
\begin{equation}
  Z(T,\mu) = \tr\!\left[e^{-(H - \mu N_q)/T}\right].
\end{equation}
We have added to the total action a term introducing a quark chemical potential
$\mu$ which multiplies the quark number operator $\hat{N}_q$. Sometimes the
baryon number operator $\hat{N}_B = \hat{N}_q/3$ and the baryon chemical
potential $\mu_B = 3\mu$ are used instead.

Since the quark number operator $\hat{N}_q$ assumes integer values, one may
expand the grand canonical partition function $Z(T,\mu)$ in a power series of
the fugacity variable
\begin{equation}
  z = e^{\mu/T}.
\end{equation}
The result of this fugacity expansion is a sum over canonical partition
functions $Z_n(T)$ with a fixed quark number $n \in \mathbb{Z}$
\begin{equation}
  Z(T,\mu) = \sum_n \left(e^{\mu/T}\right)^n Z_n(T),
\end{equation}
where negative values of $n$ correspond to a net surplus of antiquarks. On a
finite lattice the sum over the quark number $n$ is finite, i.e.,
$|n| \le n_{max}$. For a spatial volume with $N^3$ points each flavor of Wilson
fermions contributes to $n_{max}$ the number $6N^3$.

With the help of the grand canonical partition function (12.20) we can compute
new observables such as the expectation value of the quark number density
\begin{equation}
  n_q \equiv \frac{1}{V}\ev{\hat{N}_q}
    = \frac{T}{V}\frac{\partial\ln Z(T,\mu)}{\partial\mu}.
\end{equation}
The quark number density of the physical environment will affect the phase
structure and the location and properties of the finite temperature phase
transition.

In addition to the gauge coupling and the quark mass parameters, familiar from
previous chapters, and the temperature $T$ introduced in the last section, we
now have the additional parameter $\mu$. The phase diagram will become
increasingly complicated, as we will see in the discussion in Sect.~12.2.2.

\subsection{The chemical potential on the lattice}

How can we introduce the quark number density in the lattice action? The
Euclidean continuum quark number operator is given by the spatial volume
integral over the temporal component
\begin{equation}
  \bar{\psi}(x)\,\gamma_4\,\psi(x)
\end{equation}
of the conserved vector current $\bar{\psi}\gamma_\mu\psi$. In a lattice
implementation one is tempted to add to the Dirac operator (2.36) such a term
multiplied with $\mu$, since the factor $1/T$ comes from integrating over the
finite time direction. That simplistic ansatz runs into problems, though. In
order to see this let us discuss the free case. We will show that there the
energy density $\varepsilon(\mu)$ becomes divergent in the continuum limit and
that this unphysical divergence is a consequence of the naive introduction of
the chemical potential according to (12.24).

The free energy density is given by
\begin{equation}
  \varepsilon(\mu) = -\frac{1}{V_3}\ev{\hat{H}}
    = -\frac{1}{V_3\,Z}\frac{\partial Z}{\partial\beta}
    = -\frac{1}{V_3}\frac{\partial\ln Z}{\partial\beta},
\end{equation}
where $V_3$ is the spatial volume and $Z = \tr\exp(-\beta\hat{H})$. When
evaluating the energy density on the lattice it is convenient to introduce a
different lattice spacing $a_T$ for the temporal direction, which in the end of
the calculation can be set equal to the lattice spacing $a$ used for the
spatial direction. Due to the fact that for the free case the partition
function is $Z = \det[D]$, we find for the energy density (using
$\beta = N_T a_T$)
\begin{equation}
\begin{split}
  \varepsilon(\mu) &= -\frac{1}{(aN)^3}\,
    \frac{\partial}{\partial(N_T a_T)}\ln(\det[D])\Bigg|_{\mu a_T=\mathrm{const}}\\
  &= -\frac{1}{(aN)^3\,N_T}\,
    \frac{\partial}{\partial(a_T)}\ln(\det[D])\Bigg|_{\mu a_T=\mathrm{const}}.
\end{split}
\end{equation}

For the evaluation of $\varepsilon(\mu)$ we need to compute the determinant of
the lattice Dirac operator $D$. For free fermions with a chemical potential
introduced in the naive way (12.24) it reads
\begin{equation}
  D(n|m) = \sum_{j=1}^3 \gamma_j\,
    \frac{\delta_{n+\hat{j},m} - \delta_{n-\hat{j},m}}{2a}
    + \gamma_4\,\frac{\delta_{n+\hat{4},m} - \delta_{n-\hat{4},m}}{2a_T}
    + m\,\delta_{nm} + \mu\,\gamma_4\,\delta_{nm}.
\end{equation}

For the evaluation of the determinant we Fourier transform $D$ to momentum
space (cf.\ Appendix A.3) and obtain
\begin{equation}
  D(p|q) = \frac{1}{|\Lambda|}\sum_{n,m\in\Lambda} e^{-ip\cdot n}\,
    D(n|m)\,e^{iq\cdot m} = \delta(p-q)\,\tilde{D}(p),
\end{equation}
with the Dirac operator in momentum space $\tilde{D}$ given by
\begin{equation}
  \tilde{D}(p) = \frac{1}{a_T}\left[
    i\,\frac{a_T}{a}\sum_{j=1}^3\gamma_j\sin(p_j a)
    + i\gamma_4\sin(p_4 a_T) + a_T m\,\mathbb{1} + a_T\mu\,\gamma_4
  \right].
\end{equation}

For the logarithm of the determinant needed in (12.26) we find (use (A.54))
\begin{equation}
  \ln(\det[D]) = \ln\prod_p \det_d\tilde{D}(p)
    = \sum_p \ln\det_d\tilde{D}(p)
    = \sum_p \tr_d\!\left[\ln\tilde{D}(p)\right].
\end{equation}
We use the subscript $d$ attached to the determinant and the trace to indicate
that both are computed only for the remaining $4\times 4$ matrices in Dirac
space. Using (12.26), (12.27), (12.28), (12.29), and (12.30) we obtain for the
derivative with respect to $a_T$,
\begin{equation}
\begin{split}
  &\frac{\partial}{\partial a_T}\ln(\det[D])\Bigg|_{\mu a_T=\mathrm{const}}\\
  &\quad= \sum_p \tr_d\!\left[
    \frac{\partial}{\partial a_T}\ln\tilde{D}(p)
  \right]_{\mu a_T=\mathrm{const}} \\
  &\quad= C + \sum_p \tr_d\!\left[
    \left(a_T\tilde{D}(p)\right)^{-1}\left(
      \frac{i}{a}\sum_{j=1}^3\gamma_j\sin(p_j a) + m\,\mathbb{1}
    \right)
  \right],
\end{split}
\end{equation}
where the constant $C$ is given by $C = -4N_T N^3/a_T$. Evaluating the
remaining traces in Dirac space and inserting the result into (12.26) we find
after setting $a_T = a$,
\begin{equation}
  \varepsilon(\mu) = C - \frac{4}{N^3 N_T a^4}\sum_p F(ap, am, a\mu),
\end{equation}
with
\begin{equation}
  F(ap, am, a\mu)
  = \frac{\sum_{j=1}^3\sin^2(a p_j) + (am)^2}
         {\sum_{k=1}^3\sin^2(a p_k) + (am)^2
          + \left(\sin(a p_4) - i\,a\mu\right)^2}.
\end{equation}

When computing the final result for the energy density we need to normalize by
removing the value for $\mu = 0$, i.e., we consider
$\varepsilon(\mu) - \varepsilon(0)$. We return to zero temperature and, letting
$N = N_T \to \infty$, we change to continuous integration over momenta $p$:
\begin{equation}
  \frac{1}{N^3 N_T a^4}\sum_p
  \;\Rightarrow\;
  \frac{1}{(2\pi)^4}\int_{-\pi}^{\pi}\dd^4 p
\end{equation}
giving rise to
\begin{equation}
  \varepsilon(\mu) - \varepsilon(0) = -\frac{4}{(2\pi)^4}\int_{-\pi}^{\pi}
    \dd^4 p\;\big(F(ap,am,a\mu) - F(ap,am,0)\big),
\end{equation}
where in $F$ we have still kept $a$ finite. Using a contour integral in the
complex $p_4$-plane, the leading term in the limit $a\to 0$ may be extracted.
From the $p_4 \to 0$ contribution one obtains the behavior
\begin{equation}
  \lim_{a\to 0}\big(\varepsilon(\mu) - \varepsilon(0)\big) \propto
    \left(\frac{\mu}{a}\right)^2 .
\end{equation}

This is divergent in the continuum limit and thus simply adding the term
(12.24) in the action does not give rise to a proper discretization of the
chemical potential.

Closer inspection of the continuum situation (cf.\ [44--46]) clarifies the
problem. The quark number is the conserved charge of the U(1) global symmetry.
Determining the Noether current for the lattice action gives the current
expressed by nearest neighbor terms. The space integral then produces a
suitable form of the chemical potential term. One therefore implements the
chemical potential by replacing the temporal hopping term in (5.51) with
\begin{equation}
\begin{split}
  -\frac{1}{2a}\sum_{n\in\Lambda}\Big[
  &f(a\mu)\,(1-\gamma_4)_{\alpha\beta}\,U_4(n)_{ab}\,\delta_{n+\hat{4},m}\\
  &\quad - f(a\mu)^{-1}\,(1+\gamma_4)_{\alpha\beta}\,
    U_4(n-\hat{4})^\dagger_{ab}\,\delta_{n-\hat{4},m}
  \Big],
\end{split}
\end{equation}
where $f(a\mu)$ is a real function not yet specified.

For $\mu = 0$ the original action should be recovered, thus we request
$f(0) = 1$. In an expansion in $\mu$ the next term should be linear in the
chemical potential in order to reproduce the density term. We therefore have
$f(a\mu) = 1 + a\mu + \mathcal{O}(a\mu)^2$. Time reflection invariance
(cf.\ Sect.~5.4) requires $f(a\mu) = 1/f(-a\mu)$. The simplest choice
fulfilling all these conditions is
\begin{equation}
  f(a\mu) = \exp(a\mu).
\end{equation}

In this formulation the propagation forward in time is favored by a factor of
$\exp(a\mu)$, whereas propagation backward in time is disfavored by
$\exp(-a\mu)$. This introduces the desired particle--antiparticle asymmetry.

In Sect.~5.3.1 we have shown that the hopping expansion expresses the fermion
determinant as a sum over closed loops. In such a loop the forward hopping
factors $f(a\mu)$ and the backward hopping factors $f(a\mu)^{-1}$ cancel unless
the loop winds nontrivially around the compact time direction. The total
contribution of the chemical potential then is $f(a\mu)^{wN_T}$, where
$w \in \mathbb{Z}$ denotes the number of windings of the loop.

Based on that observation, we note that in a simulation instead of modifying
all link terms in time direction one may also modify just all forward
time-directed hopping terms in only a single time slice with the factor
\begin{equation}
  f(a\mu)^{N_T} = \exp(a\mu N_T) = \exp(\mu/T),
\end{equation}
and the corresponding backward-oriented terms with the inverse factor. Like in
(12.20) we find that the chemical potential always enters in the form $\mu/T$.

The introduction of the chemical potential comes with a serious technical
drawback: For $a\mu \ne 0$ the Dirac operator is no longer
$\gamma_5$-hermitian. It is easy to see that multiplying the hopping term
(12.37) with $\gamma_5$ from the left and the right changes the sign of
$\gamma_4$; hermitian conjugation then exchanges $U$ and $U^\dagger$ and we
find instead of the $\gamma_5$-hermiticity relation (5.76) the modified
equation
\begin{equation}
  \gamma_5\,D(\mu)\,\gamma_5 = D^\dagger(-\mu).
\end{equation}

This, however, is no longer a symmetry transformation but leads to a term with
$f$ replaced by $1/f^*$. We find
\begin{equation}
  \gamma_5\,D(f)\,\gamma_5 = D^\dagger(1/f^*)
  \quad\Rightarrow\quad
  \det[D(f)] = \left(\det[D(1/f^*)]\right)^* .
\end{equation}

This is an invariance operation only for $f = 1/f^*$. For real $f$ one
therefore has $f = 1/f = 1$ which implies $a\mu = 0$. Consequently for
nonvanishing real $\mu$ the determinant of the Dirac operator then is complex
and one cannot obtain a real Boltzmann weight by doubling the number of fermion
flavors. A nonvanishing real $\mu$ creates a particle--antiparticle asymmetry
which destroys the reality of the determinant and a straightforward application
of importance sampling.

It is interesting to notice, however, that for imaginary chemical potential
$\mu = i\eta$, $\eta \in \mathbb{R}$ we have
\begin{equation}
  f(ia\eta) = 1/f(ia\eta)^* = 1/f(-ia\eta).
\end{equation}

Then the Dirac operator is $\gamma_5$-hermitian in the standard way and the
determinant is real [47--49]. Based on this observation one might try to work
at imaginary chemical potential and analytically continue to real values
(cf.\ Sect.~12.3.3).

Another special case is when one considers the so-called isospin chemical
potential $\mu_I$ [49, 50]. Let us introduce a chemical potential variable for
each flavor. In the Euclidean continuum Lagrangian this corresponds to terms
\begin{equation}
  \sum_f \mu_f\,\bar{\psi}_f\,\gamma_4\,\psi_f,
\end{equation}
which may be translated to the lattice formulation as discussed above.
Focusing on two light flavors we may consider $\mu_u = \mu_I$ and
$\mu_d = -\mu_I$. The term in the Lagrangian then has the form
$\mu_I(\bar{u}\gamma_4 u - \bar{d}\gamma_4 d)$. In this case the Dirac
operators are block diagonal, where the blocks are the one-flavor operators for
the up and the down quark:
\begin{equation}
  \begin{pmatrix} D(\mu_I) & 0 \\ 0 & D(-\mu_I) \end{pmatrix}
  =
  \begin{pmatrix} D(\mu_I) & 0 \\
                  0 & \gamma_5\,D^\dagger(\mu_I)\,\gamma_5 \end{pmatrix}.
\end{equation}
In this step we have used (12.40). For mass-degenerate quarks the determinant
of this matrix is real and positive since
\begin{equation}
  \det[D(\mu_I)]\,\det[\gamma_5\,D^\dagger(\mu_I)\,\gamma_5]
  = \det[D(\mu_I)]\,\det[D^\dagger(\mu_I)]
  = \left|\det[D(\mu_I)]\right|^2 .
\end{equation}
This property allows one to study nonvanishing isospin chemical potential in
Monte Carlo simulations using the standard techniques discussed in Chap.~8.
Isospin chemical potential is not an academic problem, since one can imagine
such a situation in relativistic heavy ion collisions.

There is also a definition of the overlap Dirac operator for nonvanishing
chemical potential [51, 52]. The kernel operator may be the usual Wilson Dirac
operator with chemical potential $\mu$ introduced as in (12.37), which is then
not $\gamma_5$-hermitian. The sign function applied in the construction of
$D_{ov}$ (cf.\ Sect.~7.4) is a generalization to nonhermitian but
diagonalizable operators, i.e., in the spectral representation (7.84) one takes
the sign of the real part of the complex eigenvalues. $D_{ov}(\mu)$ is not
$\gamma_5$-hermitian for $\mu \ne 0$ but inherits property (12.40). It is no
longer normal and its eigenvalues are not restricted to the Ginsparg--Wilson
circle. However, eigenvectors $\psi$ and $\gamma_5\psi$ have still eigenvalues
related by $\lambda$ and $\lambda/(\lambda - 1)$. It was demonstrated [53, 54]
that this generalization reproduces correctly the thermodynamical properties of
a free Fermi gas. The exact manifestation of chiral symmetry on the lattice in
the presence of a nonvanishing chemical potential is still discussed, however.

\subsection{The QCD phase diagram in the $(T,\mu)$ space}

The situation concerning details of the phase structure, like the exact
position of phase transitions and values of critical exponents, is yet to be
settled. Lattice results from different kinds of Dirac operators often disagree
and the numerics and statistics may not be sufficient. The phase structure
discussed here therefore can only represent the present minimal consensus
[55, 56].

Figure 12.5 shows the situation for two massless quarks and for two light
quarks close to the physical mass.

\begin{figure}[htbp]
  \centering
  \includegraphics[width=0.9\textwidth]{fig125.png}
  \caption*{Fig. 12.5. \emph{The conjectured phase diagram in the $(T,\mu)$-plane
    for (l.h.s.) two massless quarks and (r.h.s.) two mass-degenerate light
    quarks. Lattice calculations so far have been useful only for small values
    of the chemical potential, attempting to locate the tricritical point
    (l.h.s) or second-order endpoint (r.h.s.) of the first-order phase
    transition line.}}
  \label{fig:125}
\end{figure}

The temperature transition at $\mu = 0$ has been discussed in Sect.~12.1.2: For
$\Nf = 2$ it is either second order or weakly first order. For increasing quark
mass the transition weakens and eventually becomes a crossover. By introducing
a light third dynamical quark the phase transition becomes strong and the
first-order line from the middle of the phase diagram extends toward the
$\mu = 0$ region. If the mass of the third quark increases, again the situation
is that of the right-hand side plot of Fig.~12.5.

There are no lattice results for zero temperature and nonvanishing real
chemical potential. Straightforward perturbative treatment of QCD at high
density fails as well [55]. Much of the information there comes from
NJL-inspired calculations [57], Dyson--Schwinger gap equations with various
interaction models like the instanton liquid model [58] and random matrix
models [59, 60]. In the confinement phase one first has nuclear matter with
(for two massless flavors) a small transition line extending into the phase,
which separates between a resonance gas-like and a liquid-like structure. Model
calculations then indicate a first-order transition between the nuclear matter
region for $\mu_B$ above 1~GeV and, at higher density, a (two-flavor) color
superconducting phase (2SC), where the quarks pair in a flavor singlet color
anti-triplet channel [57, 58]. Adding a strange quark with a mass close to its
physical value the region around the phase transition may be more complicated,
e.g., allowing for an intermediate pocket of a crystalline color superconducting
phase. However, all calculations agree that for larger $\mu$ one eventually
enters a color--flavor-locked (CFL) phase with a nonvanishing diquark
condensate [61].

Lattice results so far have been confined to a strip close to vanishing
chemical potential, typically $\mu/T < 1$. All methods have been based on
analytic continuation from $\mu = 0$ or imaginary $\mu$ to nonvanishing $\mu$.
Of particular importance for heavy ion collisions is whether the physical
situation is close to a crossover or phase transition line and in what phase
one ends up in various experiments.

\section{Chemical potential: Monte Carlo techniques}

The quest for the optimal simulation strategy for nonvanishing chemical
potential is far from being settled. We therefore refrain from trying to cover
all results available but just discuss the underlying concepts. All approaches
so far have analytically continued results obtained from measurements for real
determinants to the actual parameter values one is interested in. There have
been two groups of such extrapolations:
\begin{itemize}
  \item Using measurements at $\mu = 0$ and extrapolating with the help of
        reweighting or power series in $\mu/T$ (moment expansion).
  \item Using results determined for purely imaginary $\mu$ (negative $\mu^2$)
        and analytic continuation to real $\mu$ (positive $\mu^2$) via a power
        series and Pad\'e rational expansion or reconstruction of the fugacity
        expansion coefficients via Fourier transformation.
\end{itemize}

Quenched simulations [62--64] first led to confusing results. It was found that
the critical value of $\mu$ decreased with the pion mass: $\mu_c \propto
m_\pi/2$, which would vanish in the chiral limit. On the other hand at zero
temperature one expects that the transition is near $\mu_c \approx m_p/3$, since
the proton is the lightest state with nonvanishing baryon number. Stephanov
[65] has pointed out that for $\mu = 0$ the quenched theory is not the
$\Nf \to 0$ limit of QCD but of a theory with $\Nf$ quarks and $\Nf$ conjugate
quarks. This implies that for a simulation the phase of the quark determinant
is of importance and dynamical fermions are necessary to obtain
$\mu_c \approx m_p/3$. Most simulations in that context are therefore done with
dynamical fermions, mainly of the staggered type.

Large temperature $T \gg T_c$ corresponds to very small time extension and one
expects that the system approaches effectively a 3D gauge theory. The temporal
loops then play the role of an adjoint Higgs field [66, 67]. This dimensional
reduction opens the way for Monte Carlo studies of the effective theory
including chemical potential [68--71]. We restrict our discussion to methods
for the full 4D system.

\subsection{Reweighting}

Reweighting is a standard tool in Monte Carlo approaches for statistical spin
systems and has been successfully used to improve interpolation between Monte
Carlo results at different couplings and for analytical continuation from real
to complex couplings. It has been particularly useful for determining the
positions of partition function zeroes for complex couplings (Lee--Yang zeroes
and Fisher zeroes).

The strategy may be understood as a more general concept which we discuss here
briefly for a simple bosonic system. The partition function is rewritten as a
sum over the possible values of the total energy of the system:
\begin{equation}
\begin{split}
  Z(\kappa) \equiv \int \mathcal{D}[\Phi]\,e^{-\kappa S[\Phi]}
    = \int\dd E\,\rho(E)\,e^{-\kappa E},
  \qquad \text{with} \\
  \rho(E) = \int \mathcal{D}[\Phi]\,\delta(E - S[\Phi]).
\end{split}
\end{equation}

In this way the energy distribution density $\rho(E)$ may be obtained with
high precision by combining Monte Carlo results from different values of the
coupling -- the so-called multi-histogram technique [72, 73]. Once the density
of states $\rho(E)$ is known, $Z(\kappa)$ and its derivatives can be computed
for any value of the coupling $\kappa$. This is correct, however, only if the
energy distribution is known to arbitrary precision. In realistic simulations
this is not the case and therefore the technically crucial point is how much
the distribution generated from Monte Carlo runs at some couplings overlaps
with that at some other values of the coupling where one wants to compute $Z$.

In the spirit of using the determinant as a weight factor one may use the
deviation of $\det[D(\mu)]$ from $\det[D(0)]$ as a reweighting factor. From
\begin{equation}
\begin{split}
  &\left\langle
    \frac{e^{-S_G[\beta',U]}\,\det[D(\mu',U)]}
         {e^{-S_G[\beta,U]}\,\det[D(\mu,U)]}
  \right\rangle_{\beta,\mu}\\
  &\quad= \frac{1}{Z(\beta,\mu)}\int\mathcal{D}[U]\,
     \frac{e^{-S_G[\beta',U]}\,\det[D(\mu',U)]}
          {e^{-S_G[\beta,U]}\,\det[D(\mu,U)]}\,
     e^{-S_G[\beta,U]}\,\det[D(\mu,U)] \\
  &\quad= \frac{1}{Z(\beta,\mu)}\int\mathcal{D}[U]\,
     e^{-S_G[\beta',U]}\,\det[D(\mu',U)]
  = Z(\beta',\mu')/Z(\beta,\mu),
\end{split}
\end{equation}
we find that we may, at least in theory, determine $Z(\beta',\mu')$ from results
evaluated at $\beta$ and $\mu$. We therefore could obtain values for
\begin{equation}
  Z(\beta',\mu') = Z(\beta,\mu)\,
  \left\langle
    \frac{e^{-S_G[\beta',U]}\,\det[D(\mu',U)]}
         {e^{-S_G[\beta,U]}\,\det[D(\mu,U)]}
  \right\rangle_{\beta,\mu}.
\end{equation}

(In these equations $\beta$ denotes the inverse gauge coupling, not the inverse
temperature.) Since in a Monte Carlo sampling the configurations are obtained
with probabilities determined by the values of $\beta$ and $\mu$, the results
crucially depend on how far one wants to extrapolate. Studying the finite
temperature phase transition one then samples around
$(\beta = \beta_c, \mu = 0)$, where $\beta_c$ denotes the value of the gauge
coupling for which $aN_T = 1/T_c$.

The technical issue is the overlap of the ensemble sampled at $\beta,\mu$ with
the actual ensemble at the target couplings $\beta',\mu'$, and the resulting
statistical reliability. The difference in the total free energy grows with the
lattice volume rendering the extrapolation for larger lattice sizes more
problematic. With increasing chemical potential it becomes worse and the
predictive power poorer.

The first attempts have kept the gauge coupling fixed [62, 64] and extrapolated
from simulation at vanishing chemical potential to a nonvanishing value for
small lattices.

Fodor and Katz [74, 75] extrapolate by reweighting in both variables, the gauge
coupling and the chemical potential, sampling at the transition point
$(\beta = \beta_c, \mu = 0)$, following the pseudo-critical line. This strategy
appears to improve the approximation and was applied on small lattices.
Analyzing the complex (Lee--Yang) zeroes of the partition function allows one
to locate the endpoint of the expected first-order transition line. In a
simulation with 2+1-staggered quarks near their physical mass values the point
was estimated to lie at $\mu_c^E \approx 120(13)$~MeV at $T \approx 162(2)$~MeV
[76], corresponding to $\mu_c^E/T \approx 0.75$.

A systematic check of the convergence properties of the reweighting approach is
possible for imaginary chemical potential which can be simulated by standard
methods. For a critical assessment of such a comparison, cf.\ [77--79].

\subsection{Series expansion}

The chemical potential always enters in the combination $\mu/T$. In another
method of extrapolation one therefore writes the physical observables as a
power series in $\mu/T$. Time-reversal invariance of the gauge configuration
ensemble (compare Sect.~5.4) provides the symmetry
\begin{equation}
  Z(\mu/T) = Z(-\mu/T).
\end{equation}

Thus $Z$ and observables that are symmetric under time reflection are even
functions in $\mu/T$ and the expansion will be in even powers $(\mu/T)^2$. The
pressure density (12.18) in a homogeneous system is $p = (T/V)\ln Z$. The
corresponding Taylor--McLaurin series then is even in $\mu$ and reads
\begin{equation}
  p = T^4\sum_{n=0}^{\infty} c_{2n}(T)\left(\frac{\mu}{T}\right)^{2n},
\end{equation}
where the coefficients are derivatives of the free energy $F = -\ln(Z)/\beta$
and therefore generalized quark number susceptibilities, evaluated for
vanishing $\mu$ [80--83].

The coefficient of the quadratic term in (12.50) is the quark number
susceptibility at $\mu = 0$. It is a derivative of the quark number density
(12.23) and a second derivative of the fermion determinant with respect to the
chemical potential,
\begin{equation}
  \chi_q = \frac{\partial n_q}{\partial\mu}
    = \frac{T}{V}\frac{\partial^2\ln Z}{\partial\mu^2}.
\end{equation}
Such derivatives involve traces like
\begin{equation}
  \frac{\partial\ln\det[D]}{\partial\mu}
  = \tr\!\left[D^{-1}\frac{\partial D}{\partial\mu}\right]
\end{equation}
and higher derivatives (cf.\ [83]). These are evaluated in a similar way as has
been discussed in the context of the fermionic force in (8.44) and (8.47). The
trace is evaluated with the help of stochastic estimators as discussed in
Sect.~8.4. The coefficients in (12.50) are determined in simulations at
$\mu = 0$.

In the thermodynamic limit phase transitions will be associated with
singularities in some observables. The coefficients of the expansion can be
used to estimate the position of the closest singularity. The quark number
susceptibility exhibits a pronounced peak at the transition [84]. Finite size
scaling considerations then may give hints on the type of phase transition or
crossover. Several groups (in particular [82--86]) have analyzed the quark
number susceptibility and other expansion coefficients in this way.

\subsection{Imaginary $\mu$}

In (12.42) we have noted that for imaginary chemical potential $\mu = i\eta$,
$\eta \in \mathbb{R}$, the Dirac operator is $\gamma_5$-hermitian and the
determinant is real. Then standard importance sampling techniques can be
applied like for the case without chemical potential.

The chemical potential enters the partition function through the factors
$\exp(\pm i\eta/T)$ as can be seen from (12.39). Thus thermodynamic functions
are $2\pi$-periodic in $\eta/T$ and we can restrict $\eta/T$ to the interval
$[0, 2\pi)$.

Let us discuss the analytic structure in the complex $z$-plane [87] where
\begin{equation}
  z \equiv (\mu/T)^2 = -(\eta/T)^2 .
\end{equation}

Dynamical fermions break the $Z_3$ symmetry of the pure gauge theory (see
Sect.~12.1.2). However, the effect of a $Z_3$ transformation on the fermion
fields can be compensated by a shift in the complex chemical potential, thus
leading to a new symmetry. This implies a periodicity in the imaginary part of
the chemical potential,
\begin{equation}
  Z(\eta/T) = Z(\eta/T + 2\pi/3) .
\end{equation}

Thus the interval $\eta/T \in [0, 2\pi)$ is split into three equivalent sectors
distinguished only by the phase of the Polyakov loop. For imaginary chemical
potential there are $Z_3$ transitions of first order at the boundaries of the
sectors and crossover lines above a certain temperature [88, 89]. These
singularities limit the domain of convergence for the series to roughly
$|\mu/T| \le 1$.

One can directly simulate the theory at such negative $z$-values, i.e.,
imaginary chemical potential $\mu = i\eta$, and attempt to extrapolate the
results to positive $z$ [88--92]. Analytic continuation may be performed in
several ways. The simplest approach is a power series in $z$, which converges
in a circular domain $|z - z_0| < R$, where $R$ is the distance to the closest
singularity. A standard method in the theory of analytic functions is to
perform a sequence of expansions around points located suitable in the
convergence domain of the preceding series. Other methods involve optimal
mappings or Pad\'e expansions. The latter is a systematic method to replace the
power series by that rational function which has identical expansion
coefficients.

A straightforward approach is to fit the simulation results at several values
of $z$ to a power series as in (12.50). The Taylor expansion has the same
convergence properties as shown for the pressure in [81, 83]. It turns out that
convergence is rapid for several observables studied, like the chiral
condensate and screening masses [69], as well as the position of the crossover
(the pseudo-critical temperature) [79, 88, 89, 92]. Pad\'e approximants may
allow extension of the extrapolation range [93] beyond the convergence circle
of the power series.

Another extension to complex values of $\mu$ introduces variables
$\cosh(\mu a)$ and $\sinh(\mu a)$ [94, 95], with subsequent analytic
extrapolation to real chemical potential. The results for the position of the
crossover are in agreement with the above-mentioned work.

\subsection{Canonical partition functions}

For imaginary chemical potential the fugacity expansion of the grand canonical
partition function (12.22) is a Fourier sum [47, 48] and periodic in $\eta/T$
with a period of $2\pi$,
\begin{equation}
  Z(T, i\eta) = \sum_n e^{i n\eta/T}\,Z_n(T).
\end{equation}

This sum may be inverted leading to the following expression for the canonical
partition functions $Z_n(T)$ at a fixed quark number $n$:
\begin{equation}
  Z_n(T) = \frac{1}{2\pi}\int_{-\pi}^{\pi}\dd(\eta/T)\;
    e^{-i n\eta/T}\,Z(T, i\eta).
\end{equation}
Equation (12.56) thus is an exact expression for the coefficients $Z_n$ of the
fugacity expansion (12.22).

Simulations at imaginary chemical potential can be used to determine values for
$Z(T,i\eta)/Z(T,0)$ and thus in principle the Fourier integral (12.56) for the
$Z_n$ can be evaluated numerically. A problem is the oscillations of the Fourier
sum which increase with the quark number $n$, such that at higher values of
$|n|$ many intermediate points of $\eta/T$ in the interval $[-\pi,\pi]$ are
necessary for a reliable numerical estimate of the Fourier integral.
Combinations of multi-histogram sampling and other intricate methods can be
used to determine $Z_n$ and then, via (12.22), $Z(T,\mu)$ [78, 96--101]. An
alternative approach that has been explored recently [102, 103] is perturbative
techniques for the evaluation of the $Z_n$.

Studying the canonical partition function is particularly suitable to explore
few nucleon systems at low temperature, and in principle, allows one to study
the bulk properties of nuclear matter and the nuclear interactions. The
metastable region is nicely exhibited and a Maxwell construction relates this
signal to the grand canonical results. Nice agreement of these results from
this approach with other approaches discussed earlier has been demonstrated
[98]. The results so far have been on small lattices. Fixing the baryon
density, though, requires to increase $n$ for increasing volumes. This may be a
problem for the necessary statistics of the simulations.

% [unrecovered: the reference list in extract/ch12.txt is truncated after
% entry 52; entries 53-103, which are cited in the text, were not extracted]

\section*{References}

\begin{enumerate}[label={[\arabic*]}]
  \item Z.~Fodor: PoS LATTICE2007, 011 (2007) 301, 308.
  \item F.~Karsch: PoS LATTICE2007, 015 (2007) 301, 308.
  \item O.~Philipsen: in \emph{Conceptual and Numerical Challenges in Femto-
        and Peta-Scale Physics}, edited by C.~Gattringer et al., The European
        Physical Journal ST 152, p.~29 (Springer, Berlin, Heidelberg, New York
        2007) 301.
  \item C.~DeTar: PoS LATTICE2008, 001 (2008) 301.
  \item S.~Ejiri: PoS LATTICE2008, 002 (2008) 301.
  \item A.~M.~Polyakov: Phys.\ Lett.\ B 72, 477 (1978) 303.
  \item L.~Susskind: Phys.\ Rev.\ D 20, 2610 (1979) 303.
  \item G.~'t Hooft: Nucl.\ Phys.\ B 153, 141 (1979) 305.
  \item C.~B.~Lang and M.~Wiltgen: Phys.\ Lett.\ B 131, 153 (1983) 305.
  \item G.~Bali: Phys.\ Rep.\ 343, 1 (2001) 305.
  \item H.~Ichie, V.~Bornyakov, T.~Streuer, and G.~Schierholz: Nucl.\ Phys.\ A
        721, 899 (2003) 305.
  \item F.~Bissey et al.: Phys.\ Rev.\ D 76, 114512 (2007) 305.
  \item Y.~Peng and R.~W.~Haymaker: Phys.\ Rev.\ D 47, 5104 (1993) 305.
  \item B.~Svetitsky and L.~G.~Yaffe: Nucl.\ Phys.\ B 210, 423 (1982)
        305, 306.
  \item B.~Svetitsky and L.~G.~Yaffe: Phys.\ Rev.\ D 26, 963 (1982) 305.
  \item L.~D.~McLerran and B.~Svetitsky: Phys.\ Lett.\ B 98, 195 (1981) 306.
  \item J.~Kuti, J.~Polonyi, and K.~Szlachanyi: Phys.\ Lett.\ B 98, 199 (1981)
        306.
  \item L.~D.~McLerran and B.~Svetitsky: Phys.\ Rev.\ D 24, 450 (1981) 306.
  \item E.~Laermann and O.~Philipsen: Ann.\ Rev.\ Nucl.\ Part.\ Sci.\ 53, 163
        (2003) 306.
  \item F.~Karsch and E.~Laermann: in \emph{Quark Gluon Plasma}, edited by\\
        R.~C.~Hwa, p.~1 (World Scientific, Singapore 2003) 306.
  \item J.~L.~Cardy: \emph{Finite-Size Scaling} (North-Holland, Amsterdam
        1988) 306.
  \item J.~Engels, F.~Karsch, I.~Montvay, and H.~Satz: Phys.\ Lett.\ B 101, 89
        (1981) 307.
  \item J.~Engels, F.~Karsch, I.~Montvay, and H.~Satz: Nucl.\ Phys.\ B 205
        [FS5], 545 (1982) 307.
  \item I.~Montvay and G.~M\"unster: \emph{Quantum Fields on a Lattice}
        (Cambridge University Press, Cambridge, New York 1994) 307.
  \item K.~Fredenhagen and M.~Marcu: Commun.\ Math.\ Phys.\ 92, 81 (1983) 307.
  \item J.~Bricmont and J.~Fr\"ohlich: Phys.\ Lett.\ B 122, 73 (1983) 307.
  \item Y.~Aoki, Z.~Fodor, S.~D.~Katz, and K.~K.~Szabo: Phys.\ Lett.\ B 643, 46
        (2006) 308.
  \item M.~Cheng et al.: Phys.\ Rev.\ D 74, 054507 (2006) 308.
  \item F.~Karsch: Lect.\ Notes Phys.\ 583, 209 (2002) 309.
  \item Y.~Aoki, Z.~Fodor, S.~D.~Katz, and K.~K.~Szabo: JHEP 0601, 089 (2006)
        309.
  \item C.~Bernard et al.: PoS LAT2005, 156 (2005) 309.
  \item R.~D.~Pisarski and F.~Wilczek: Phys.\ Rev.\ D 29, 338 (1984) 309.
  \item E.~Shuryak: Phys.\ Rep.\ 61, 71 (1980) 310.
  \item R.~P.~Feynman: \emph{Statistical Mechanics} (Westview Press, Boulder
        1998) 310.
  \item F.~Karsch, E.~Laermann, and A.~Peikert: Phys.\ Lett.\ B 478, 447 (2000)
        311.
  \item F.~Karsch: in \emph{AIP Conf.\ Proc.\ (PANIC05)}, Vol.~842, p.~20
        (Am.\ Inst.\ Phys., Melville, New York 2006) 311.
  \item P.~de Forcrand et al.: Phys.\ Rev.\ D 63, 054501 (2001) 311.
  \item I.~Wetzorke et al.: Nucl.\ Phys.\ B (Proc.\ Suppl.) 106, 510 (2002)
        311.
  \item M.~Asakawa, T.~Hatsuda, and Y.~Nakahara: Nucl.\ Phys.\ A 715, 863
        (2003) 311.
  \item J.~B.~Kogut, J.~F.~Lagae, and D.~K.~Sinclair: Phys.\ Rev.\ D 58,
        054504 (1998) 311.
  \item I.~Pushkina et al.: Phys.\ Lett.\ B 609, 265 (2005) 311.
  \item S.~Wissel et al.: PoS LAT2005, 164 (2005) 311.
  \item R.~V.~Gavai, S.~Gupta, and R.~Lacaze: PoS LAT2006, 135 (2006) 311.
  \item P.~Hasenfratz and F.~Karsch: Phys.\ Lett.\ B 125, 308 (1983) 315.
  \item J.~Kogut, H.~Matsuoka, M.~Stone, and H.~W.~Wyld: Nucl.\ Phys.\ B 225
        [FS9], 93 (1983) 315.
  \item N.~Bilic and R.~V.~Gavai: Z.\ Phys.\ C 23, 77 (1984) 315.
  \item E.~Dagotto, A.~Moreo, R.~L.~Sugar, and D.~Toussaint: Phys.\ Rev.\ B 41,
        811 (1990) 316, 322.
  \item A.~Hasenfratz and D.~Toussaint: Nucl.\ Phys.\ B 371, 539 (1992)
        316, 322.
  \item M.~G.~Alford, A.~Kapustin, and F.~Wilczek: Phys.\ Rev.\ D 59, 054502
        (1999) 316.
  \item D.~T.~Son and M.~A.~Stephanov: Phys.\ Rev.\ Lett.\ 86, 592 (2001) 316.
  \item J.~Bloch and T.~Wettig: Phys.\ Rev.\ Lett.\ 97, 012003 (2006) 317.
  \item J.~Bloch and T.~Wettig: Phys.\ Rev.\ D 76, 114511 (2007) 317.
\end{enumerate}
